\chapter{Proof of the barrier theorem}\label{sec:barrier-proof}

In this section, we prove \cref{thm:barrier-730} and \cref{b:cor:subcritical}, and, in \cref{sec:bp-strong}, the quantitative and almost-sure forms of \cref{thm:barrier-strong}. We first make the thresholds of \cref{sec:induction} random and show, using \cref{cor:ind-witness}, that a path avoiding $V$ keeps a positive overlap with its first reference during the first $T_n=\lceil n^2e^{L}\rceil$ update attempts (\cref{prop:bp-obstruction,lem:bp-coefficient}). This conclusion is deterministic: it holds for every path that avoids $V$, and the thresholds are only a device in its proof. Combining it with stationarity, spin-flip symmetry, and a Poisson clock, we obtain bounds on the Gibbs mass of the initial states from which the dynamics is close to equilibrium (\cref{prop:bp-quenched}). Integration over the disorder event of \cref{sec:references} then proves \cref{thm:barrier-730}.

The main inputs are \cref{lem:ind-predictable}, \cref{cor:ind-witness}, the families of reference configurations of \cref{sec:references} with the bound $\pi(V)\le3e^{-4L}$ (see \eqref{b:dyn:fixed-bank}), and the properties of the heat-bath dynamics in \cref{lem:model-dynamics}. The probability of the disorder event of \cref{sec:references} enters only once, in the last step of the proof of \cref{thm:barrier-730}. The quantities that are multiplied by the number of time indices are the Gibbs mass $\pi(V)$ and probabilities in an auxiliary space of thresholds, both at a fixed disorder.

The argument follows \cite[Sections~8.1 and~8.5]{oai-barriers}. The product bound for predictable thresholds is the elimination argument following \cite[(8.28)]{oai-barriers}. The reduction of total variation to a sign passage, with the half-space $\{y:R(v,y)\le0\}$ of Gibbs mass at least $1/2$ and a Poisson clock, is \cite[(8.1)--(8.3)]{oai-barriers}, together with the remark on discrete time at the end of \cite[Section~8.5]{oai-barriers}. There are three changes. First, the windows of \cref{cor:ind-witness} have radius of order $\varepsilon$ and sit at deterministic positions. We therefore count only the $M$ windows at the first levels of the blocks, in place of the union over subsets of size $N_{\rm lev}/B$ in \cite[(8.29)]{oai-barriers}. Second, the static estimates of \cref{sec:static} are uniform over the compact interval $J_\beta$, and we need no random event restricting the thresholds, such as \cite[(8.12)]{oai-barriers}. As a result, the sign obstruction holds for every path avoiding $V$, and not only with high probability. Third, the budget $T_n$ of update attempts has logarithm $L+O(\log n)$, in place of $2L$, and every fixed $\tau<1/2$ is allowed, in place of $\tau=1/4$.

We keep the notation and the standing assumptions of \cref{sec:ind-setting}. Recall that $\rho=Rn^{-\lambda}$, $\varepsilon=E\rho$ with $E=1/100$, $\zeta=2\eta_n$, the integers $M$ and $N_{\rm lev}=MB$ and the intervals $I_i$ are given by \eqref{b:dyn:threshold-intervals}, and the radius $\Psi_0=2\varepsilon+\zeta+2/n$ is given by \eqref{b:dyn:radii}.

\section{Predictable thresholds}\label{b:sec:dyn-count}

Fix the disorder, a family of reference configurations, and a path $x_0,\dots,x_m$. On the product $I_0\times\cdots\times I_{N_{\rm lev}-1}$, let $\PP^{\rm thr}$ be the product of the uniform probability measures on the intervals. Under $\PP^{\rm thr}$, the coordinates $r_0,\dots,r_{N_{\rm lev}-1}$ of a threshold array are independent and $r_i$ is uniform on $I_i$. This is an auxiliary probability space; it is not related to the randomness of the dynamics. By \cref{lem:ind-predictable}, for every fixed time $k\le m$ the overlap
\[
 z_i(k)=R\bigl(v_i,x_k\bigr)
\]
of the reference of level $i$ with $x_k$ is a Borel function of $(r_0,\dots,r_{i-1})$. The next lemma is the elimination argument of \cite[Section~8.5]{oai-barriers}.

\begin{lemma}[Windows with predictable centers]\label{lem:bp-product}
Let $I_0,\dots,I_{\bar N-1}$ be intervals of length $h>0$, and let $r_0,\dots,r_{\bar N-1}$ be independent random variables with $r_i$ uniform on $I_i$. Let $D\subseteq\{0,\dots,\bar N-1\}$, and for $i\in D$ let $a_i\ge0$ and let $z_i$ be a Borel function of $(r_0,\dots,r_{i-1})$ (a constant if $i=0$). Then
\[
 \PP\bigl\{|z_i-r_i|\le a_i\text{ for every }i\in D\bigr\}\le\prod_{i\in D}\frac{2a_i}{h}.
\]
\end{lemma}

\begin{proof}
We argue by induction on the number of elements of $D$; the case $D=\emptyset$ is trivial. Let $i_*$ be the largest element of $D$, let $W$ be the event $\{|z_i-r_i|\le a_i$ for every $i\in D\setminus\{i_*\}\}$, and let $\mathscr H$ be the $\sigma$-algebra generated by $r_0,\dots,r_{i_*-1}$. Then $W$ and $z_{i_*}$ are $\mathscr H$-measurable, and $r_{i_*}$ is independent of $\mathscr H$. It follows that
\[
 \PP\bigl\{|z_{i_*}-r_{i_*}|\le a_{i_*}\bigm|\mathscr H\bigr\}
 =\frac{\bigl|[z_{i_*}-a_{i_*},z_{i_*}+a_{i_*}]\cap I_{i_*}\bigr|}{h}\le\frac{2a_{i_*}}{h},
\]
where $|\cdot|$ denotes length. Multiplying by $\mathbf 1_W$ and taking expectations, we obtain
\[
 \PP\bigl(W\cap\{|z_{i_*}-r_{i_*}|\le a_{i_*}\}\bigr)\le\frac{2a_{i_*}}{h}\,\PP(W),
\]
and the induction hypothesis applied to $D\setminus\{i_*\}$ completes the proof.
\end{proof}

The order of integration matters: the largest index is removed first, because the later centers may depend on the earlier thresholds. In the lemma, we condition only on earlier thresholds. We do not condition on the event that the algorithm succeeds, nor on the sets $\mathcal R_j$ of retained levels, which depend on all thresholds. A bound $2a_i/h$ larger than one is allowed.

\section{A deterministic sign obstruction}\label{sec:bp-obstruction}

We define
\begin{equation}\label{b:dyn:window-weights}
 \gamma_n=\frac{2\Psi_0}{\rho}=4E+\frac{2\zeta}{\rho}+\frac4{n\rho},
\end{equation}
which is the ratio of the length of a window of radius $\Psi_0$ to the length $\rho$ of a threshold interval. By \eqref{eq:ind-scales}, we have $\gamma_n\to4(E+E')\le11/250$, and $\gamma_n\to4E=1/25$ when $E'=0$. For integers $m\ge0$ we set
\begin{equation}\label{b:dyn:prefix-count}
 C_n(m)=(m+1)\gamma_n^M.
\end{equation}

\begin{proposition}[Sign obstruction for paths avoiding $V$]\label{prop:bp-obstruction}
Assume that \eqref{b:dyn:finite-margins} and the inclusion of the neighborhood in $J_\beta$ hold, and that the family of reference configurations satisfies (P1)--(P4). Let $m\ge0$ be an integer with $C_n(m)<1$. Then
\begin{equation}\label{b:dyn:deterministic-obstruction}
 R\bigl(v(x_0),x_m\bigr)>0\qquad\text{for every path }x_0,\dots,x_m\text{ that avoids }V.
\end{equation}
\end{proposition}

\begin{proof}
Fix a path $x_0,\dots,x_m$ that avoids $V$, and suppose that $R(v(x_0),x_m)\le0$. Run the algorithm of \cref{def:ind-algorithm} with an arbitrary threshold array $\mathbf r$. Recall that the reference $v(x_0)$ of level $0$ does not depend on $\mathbf r$. The hypothesis of \cref{prop:ind-induction} therefore holds for every threshold array. By \cref{cor:ind-witness}, every threshold array belongs to the union over $0\le k\le m$ of the events
\[
 W_k=\bigl\{\mathbf r:|z_{\ell B}(k)-r_{\ell B}|\le\Psi_0\text{ for every }\ell\in\{0,\dots,M-1\}\bigr\},
\]
where $z_i(k)=R(v_i,x_k)$ and $v_i=v_i(r_0,\dots,r_{i-1})$ is the reference of level $i$ produced by the algorithm. By \cref{lem:ind-predictable}, each $z_{\ell B}(k)$ is a Borel function of $(r_0,\dots,r_{\ell B-1})$. Applying \cref{lem:bp-product} under $\PP^{\rm thr}$, with $h=\rho$, $D=\{0,B,\dots,(M-1)B\}$, and $a_i=\Psi_0$, we obtain $\PP^{\rm thr}(W_k)\le\gamma_n^M$. Summing over $k$, we find
\[
 1=\PP^{\rm thr}\Bigl(\bigcup_{k=0}^mW_k\Bigr)\le(m+1)\gamma_n^M=C_n(m)<1,
\]
which is a contradiction.
\end{proof}

The events $W_k$ use deterministic positions and deterministic radii. Their union contains the event produced by the algorithm, but we do not assert that the algorithm produces any particular configuration of retained references, and we do not condition on such an assertion. The proof also uses no property of the crossing times other than those established in \cref{sec:induction}; in particular, no distributional statement about the path at a crossing time is involved.

\section{The choice of the coefficient}\label{sec:bp-coefficient}

The count $C_n(m)$ is small when $\log(m+1)$ is smaller than $M\log(1/\gamma_n)$. Since $\gamma_n\to4(E+E')<1$ and $M$ is of order $n^\lambda$, this allows $\log m$ of order $n^\lambda$. Fix a coefficient $A$ with
\begin{equation}\label{b:dyn:coefficient}
 0<A<\frac{t_0\log\bigl(1/(4E+4E')\bigr)}{3BR};
\end{equation}
when $E'=0$, as in the proof of \cref{thm:barrier-730}, the right side is $t_0\log25/(3BR)$. We set
\begin{equation}\label{eq:bp-horizon}
 L=L_n=An^{\lambda},\qquad T_n=\bigl\lceil n^2e^{L}\bigr\rceil.
\end{equation}
Here $L_n$ is the logarithmic time scale and $e^{L_n}$ is the time horizon. The integer $T_n$ is a budget of update attempts: the continuous-time dynamics makes on average $ne^{L_n}$ attempts up to time $e^{L_n}$, and $T_n$ exceeds this number by a factor $n$.

\begin{lemma}[The count for paths of exponential length]\label{lem:bp-coefficient}
Assume \eqref{eq:ind-scales}, and let $A$ satisfy \eqref{b:dyn:coefficient}. There exists $n_0$, depending only on $A$, $t_0$, $B$, $R$, $E$, $E'$, and the sequences $(\rho_n)$ and $(\eta_n)$, such that $C_n(m)<1$ for every $n\ge n_0$ and every integer $0\le m\le n^3e^{L}$. In particular, this holds for every $0\le m\le T_n$.
\end{lemma}

\begin{proof}
By \eqref{b:dyn:coefficient}, there exists $\gamma'\in(4E+4E',1)$ with $A<t_0\log(1/\gamma')/(3BR)$; we fix such a $\gamma'$. Since $\gamma_n\to4E+4E'$, we have $\gamma_n\le\gamma'$ for all sufficiently large $n$. By \eqref{b:dyn:threshold-intervals}, we have $M\ge t_0/(3B\rho)-1=t_0n^\lambda/(3BR)-1$. Since $C_n(m)$ is increasing in $m$, it suffices to consider the largest integer $m$ with $m\le n^3e^L$, for which $m+1\le2n^3e^{L}$. We find
\[
 \begin{aligned}
 \log C_n(m)&\le\log2+3\log n+An^\lambda-\Bigl(\frac{t_0n^\lambda}{3BR}-1\Bigr)\log\frac1{\gamma'}\\
 &=-\Bigl(\frac{t_0\log(1/\gamma')}{3BR}-A\Bigr)n^\lambda+3\log n+\log2+\log\frac1{\gamma'}.
 \end{aligned}
\]
The coefficient of $n^\lambda$ is negative and $\lambda>0$, so the right side tends to $-\infty$. Finally, $T_n\le n^2e^L+1\le n^3e^L$ because $n\ge2$.
\end{proof}

We make no attempt to optimize the admissible coefficient in \eqref{b:dyn:coefficient}; only its positivity matters for \cref{thm:barrier-730}.

\section{Stationarity, the clocks, and total variation}\label{b:sec:dyn-TV}

The following lemma turns a sign obstruction for paths into a bound on the Gibbs mass of the initial states from which the dynamics is close to equilibrium. It abstracts the sign-passage reduction of \cite[Section~8.1, (8.1)--(8.3)]{oai-barriers}. It is stated for a general kernel and uses only stationarity at deterministic times, spin-flip symmetry, and the Poisson representation of the continuous-time semigroup; for the latter representation, see \cite[Section~20.1, (20.5)--(20.6)]{levin-peres-2017}.

\begin{lemma}[From a sign obstruction to total variation]\label{lem:bp-sign-test}
Let $Q$ be a Markov kernel on $\Sigma_n$ with stationary law $\pi$, and assume that $\pi(-x)=\pi(x)$ for every $x\in\Sigma_n$. Let $v\colon\Sigma_n\to\Sigma_n$ be a map, let $V\subseteq\Sigma_n$, and let $m\ge0$ be an integer. Assume that every sequence $x_0,\dots,x_k$ with $k\le m$, with $Q(x_{a-1},x_a)>0$ for $1\le a\le k$, and with $x_a\notin V$ for $0\le a\le k$, satisfies $R(v(x_0),x_k)>0$. Then for every $\tau\in[0,1/2)$ and every integer $0\le k\le m$,
\begin{equation}\label{eq:bp-sign-discrete}
 \pi\bigl\{x:\norm{Q^k(x,\cdot)-\pi}_\TV\le\tau\bigr\}\le\frac{(k+1)\pi(V)}{1/2-\tau}.
\end{equation}
Moreover, for every $t\ge0$, the kernel $P_t=e^{nt(Q-I)}$ satisfies
\begin{equation}\label{eq:bp-sign-continuous}
 \pi\bigl\{x:\norm{P_t(x,\cdot)-\pi}_\TV\le\tau\bigr\}\le\frac{(m+1)\pi(V)+\PP\{\mathsf N>m\}}{1/2-\tau},
\end{equation}
where $\mathsf N$ has the Poisson law with mean $nt$.
\end{lemma}

\begin{proof}
For $x\in\Sigma_n$, let $G_x=\{y\in\Sigma_n:R(v(x),y)\le0\}$. The map $y\mapsto-y$ preserves $\pi$ and exchanges the sets $\{y:R(v(x),y)<0\}$ and $\{y:R(v(x),y)>0\}$. These two sets therefore have the same Gibbs mass, and
\[
 \pi(G_x)=\frac{1+\pi\{y:R(v(x),y)=0\}}2\ge\frac12.
\]
Let $\nu_x$ be a probability measure on $\Sigma_n$ for each $x$. If $\norm{\nu_x-\pi}_\TV\le\tau$, then $\nu_x(G_x)\ge\pi(G_x)-\tau\ge1/2-\tau$. Multiplying by $\pi(x)$ and summing over $x$, we obtain
\begin{equation}\label{eq:bp-half-space}
 \Bigl(\frac12-\tau\Bigr)\,\pi\bigl\{x:\norm{\nu_x-\pi}_\TV\le\tau\bigr\}\le\sum_x\pi(x)\,\nu_x(G_x).
\end{equation}
The half-space $G_x$ is determined by the initial state $x$ alone. The bound \eqref{eq:bp-half-space} is a statement about the law of the dynamics started at a fixed state, averaged over that state afterward.

Let $(X_a)_{a\ge0}$ be the Markov chain with kernel $Q$ and $X_0\sim\pi$, and set $g_k=\sum_x\pi(x)Q^k(x,G_x)=\PP\{R(v(X_0),X_k)\le0\}$. Almost surely, $Q(X_{a-1},X_a)>0$ for every $a\ge1$. For $k\le m$, the hypothesis then implies that $\{R(v(X_0),X_k)\le0\}\subseteq\{X_a\in V\text{ for some }0\le a\le k\}$, up to a null set. Since $\pi Q=\pi$, the law of $X_a$ is $\pi$ for every deterministic $a$, and by a union bound we have
\begin{equation}\label{b:dyn:stationary-sign}
 g_k\le\sum_{a=0}^k\PP\{X_a\in V\}=(k+1)\pi(V)\qquad(0\le k\le m).
\end{equation}
Together with \eqref{eq:bp-half-space} for $\nu_x=Q^k(x,\cdot)$, this proves \eqref{eq:bp-sign-discrete}. The matrices $ntQ$ and $ntI$ commute, which gives $P_t=\sum_{k\ge0}\PP\{\mathsf N=k\}Q^k$ and
\[
 \sum_x\pi(x)P_t(x,G_x)=\sum_{k\ge0}\PP\{\mathsf N=k\}\,g_k\le(m+1)\pi(V)+\PP\{\mathsf N>m\},
\]
where we used \eqref{b:dyn:stationary-sign} for $k\le m$ and $g_k\le1$ for $k>m$. With \eqref{eq:bp-half-space} for $\nu_x=P_t(x,\cdot)$, this proves \eqref{eq:bp-sign-continuous}.
\end{proof}

The union bound \eqref{b:dyn:stationary-sign} uses stationarity only at the deterministic times $0,\dots,k$, and it does not use independence along the chain. The set $V$ and the map $v$ are fixed before the chain is run.

We now apply \cref{lem:bp-sign-test} to the heat-bath kernel at a fixed disorder. Recall the time scale $L$ and the budget $T_n$ of \eqref{eq:bp-horizon}.

\begin{proposition}[Gibbs mass of initial states close to equilibrium]\label{prop:bp-quenched}
Assume \eqref{eq:ind-scales}, let $A$ satisfy \eqref{b:dyn:coefficient}, and let $L=An^\lambda$ and $T_n=\lceil n^2e^L\rceil$. There exists $n_0$, depending only on the fixed parameters, such that the following holds for every $n\ge n_0$. Let $J$ be a disorder for which there exists a family of reference configurations with properties (P1)--(P4) and with $\pi^J(V)\le3e^{-4L}$. Then for every $\tau\in[0,1/2)$,
\begin{align}
 \pi^J\bigl\{x:d^J(x,e^{L})\le\tau\bigr\}
 &\le\frac{3(T_n+1)e^{-4L}+1/n}{1/2-\tau},\label{b:dyn:continuous-mixed}\\
 \pi^J\bigl\{x:d^J_{\rm disc}(x,\lfloor e^{L}\rfloor)\le\tau\bigr\}
 &\le\frac{3(T_n+1)e^{-4L}}{1/2-\tau}.\label{b:dyn:discrete-mixed}
\end{align}
\end{proposition}

\begin{proof}
Let $n_0$ be so large that the conclusions of \cref{lem:ind-margins} and \cref{lem:bp-coefficient} hold for $n\ge n_0$, and fix a family as in the statement. We apply \cref{lem:bp-sign-test} with $Q=K^J$, the Gibbs measure $\pi=\pi^J$, the map $v(x)$ of \cref{def:ind-algorithm}, the set $V$ of the family, and $m=T_n$. The kernel $K^J$ preserves $\pi^J$ by \cref{lem:model-dynamics}(a), and $\pi^J(-x)=\pi^J(x)$ because $H_n(-x)=H_n(x)$. If $K^J(x_{a-1},x_a)>0$, then $x_a$ and $x_{a-1}$ differ in at most one coordinate by \eqref{b:eq:discrete-kernel}. Every sequence in the hypothesis of \cref{lem:bp-sign-test} is a path in the sense of \eqref{eq:ind-jump}. For $k\le T_n$, we have $C_n(k)<1$ by \cref{lem:bp-coefficient}, and \cref{prop:bp-obstruction}, applied to the path $x_0,\dots,x_k$, implies that $R(v(x_0),x_k)>0$ when the path avoids $V$. We have verified the hypothesis of \cref{lem:bp-sign-test}.

In continuous time, we have $P^J_t=e^{nt(K^J-I)}$ by \eqref{eq:common-semigroup}, and we take $t=e^L$. Then $\mathsf N$ has mean $ne^L$, and Markov's inequality gives
\begin{equation}\label{b:dyn:clock-tail}
 \PP\{\mathsf N>T_n\}\le\frac{ne^L}{T_n}\le\frac{ne^L}{n^2e^L}=\frac1n.
\end{equation}
Inserting \eqref{b:dyn:clock-tail} and $\pi^J(V)\le3e^{-4L}$ into \eqref{eq:bp-sign-continuous}, we obtain \eqref{b:dyn:continuous-mixed}. In discrete time, $k=\lfloor e^L\rfloor\le T_n$, and by \eqref{eq:bp-sign-discrete} we have
\[
 \pi^J\bigl\{x:d^J_{\rm disc}(x,\lfloor e^L\rfloor)\le\tau\bigr\}\le\frac{(\lfloor e^L\rfloor+1)\,3e^{-4L}}{1/2-\tau},
\]
which implies \eqref{b:dyn:discrete-mixed}.
\end{proof}

For every fixed $A>0$, the right sides of \eqref{b:dyn:continuous-mixed} and \eqref{b:dyn:discrete-mixed} tend to zero. Indeed, $T_n+1\le n^2e^L+2$, and we have
\begin{equation}\label{eq:bp-mixed-limit}
 3(T_n+1)e^{-4L}\le3n^2e^{-3An^\lambda}+6e^{-4An^\lambda}\longrightarrow0.
\end{equation}

\section{The parameters and the exponents}\label{b:sec:completion}

We first list the parameters of the proof of \cref{thm:barrier-730} in the order in which they are chosen. All of them are fixed before $n\to\infty$ and before the threshold $\tau$.

\begin{table}[ht]
\centering
\small
\begin{tabular}{@{}c>{\raggedright\arraybackslash}p{0.47\textwidth}>{\raggedright\arraybackslash}p{0.33\textwidth}@{}}
\toprule
Step & Parameter & Condition and reference \\
\midrule
1 & the parameters of \cref{tab:ref-parameters} up to $H$: $q_\beta$, $q_*$, $b_1$, $b_2$, $t_0$, $J_\beta$, $I_\beta$, $m_\beta$, $\vartheta$, $E$, $g$, $\varkappa$, $C_b$, $R$, the constants of the static estimates, and $H$ & \cref{sec:ref-parameters}; $C_b$ from \cref{inp:hb-rate} \\
2 & the threshold geometry: intervals of length $\rho_n$ and spacing $3\rho_n$ & \eqref{b:dyn:threshold-intervals} \\
3 & the block length $B$ & any integer $B\ge H+2$, \eqref{b:dyn:BH} \\
4 & the constants $c_s$, $C_s$ of the list sizes & \eqref{eq:ref-schedule} \\
5 & the coefficient $A_\beta$ & $0<A_\beta<t_0\log25/(3BR)$, \eqref{b:dyn:coefficient} \\
6 & for each $n$: $\rho_n$, $\varepsilon_n$, $\eta_n$, $L_n=A_\beta n^{7/30}$, $M$, $N_{\rm lev}$, the list sizes $K_s$, and $T_n$ & \eqref{b:dyn:scales}, \eqref{b:dyn:threshold-intervals}, \eqref{eq:ref-schedule}, \eqref{eq:bp-horizon} \\
7 & the threshold $\tau\in(0,1/2)$ & enters only through $1/(1/2-\tau)$ \\
\bottomrule
\end{tabular}
\caption{The parameters in the proof of \cref{thm:barrier-730}, in the order in which they are chosen. Each parameter of Steps 1--5 depends only on $\beta$ and on the parameters chosen before it.}
\label{tab:bp-order}
\end{table}

We use the exponents $\lambda=7/30$ and $\xi=5/21$. \Cref{tab:bp-exponents} lists every condition on the exponents used in the proof, with its value at these exponents; the corresponding powers of $n$ are listed in \cref{tab:ref-exponents}.

\begin{table}[ht]
\centering
\small
\begin{tabular}{@{}>{\raggedright\arraybackslash}p{0.64\textwidth}c@{}}
\toprule
Condition (where it is used) & Value at $\lambda=7/30$ \\
\midrule
$3\lambda\le a$, where $Cn^{-a}$ bounds the mean bias (denominator margin, \cref{sec:static}) & $a=7/10=3\lambda$ \\
$7/15-2+6\lambda<0$ (Chebyshev bound for the denominator) & $-2/15$ \\
$1-4\lambda>0$ (local and separated-overlap numerators, $e^{-cn\rho_n^4}$) & $1/15$ \\
$1-4\xi>0$ (fixed-cutoff numerators, $e^{-cn\eta_n^4}$) & $1/21$ \\
$3\lambda-2\xi>0$ (so that $\rho_n^3=o(\eta_n^2)$) & $47/210$ \\
$1-3\lambda>\lambda$ (conditional narrow mass $e^{-cn\rho_n^3}$ against $e^{CL_n}$) & $3/10>7/30$ \\
$1-2\xi>\lambda$ (conditional fixed-cutoff mass against $e^{CL_n}$) & $11/21>7/30$ \\
$\lambda<1$ (overlap-neighborhood bound needs $L_n=o(n)$; $1/(n\rho_n)\to0$) & yes \\
$\xi-\lambda>0$ (so that $\eta_n=o(\rho_n)$, \cref{sec:induction}) & $1/210$ \\
\bottomrule
\end{tabular}
\caption{The conditions on the exponents, with $\xi=5/21$. The entries in the last column are the values of the left sides at $\lambda=7/30$; for instance, the probability that the denominator bound fails is $O(n^{7/15-2+6\lambda})=O(n^{-2/15})$, and the probability that the bound for a single local numerator fails is at most $e^{-cn^{1-4\lambda}}=e^{-cn^{1/15}}$. In the first row, the mean-bias bound of \cref{inp:hb-rate} has $a=3\lambda$, and the strict margin in \eqref{eq:st-S1} comes from the choice of $R$ in \eqref{b:dyn:R-condition}.}
\label{tab:bp-exponents}
\end{table}

The conditions of \cref{tab:bp-exponents} also explain why the argument gives no exponent of $1/4$ or more. For an exponent $a$ of the mean-bias bound and an exponent $\lambda\le a/3$ with $\lambda<1/4$, all rows hold if $\xi$ is chosen strictly between $\lambda$ and $\min\{1/4,3\lambda/2\}$, which is the condition \eqref{eq:ref-exponents}; the choice $\xi=5/21$ works for $\lambda=7/30$. At $\lambda=1/4$, the third row fails, since then $n\rho_n^4=R^4$ does not tend to infinity.

\section{Proof of \texorpdfstring{\cref{thm:barrier-730}}{the barrier theorem}}\label{sec:bp-730}

\Cref{sec:induction} and \cref{sec:bp-obstruction,sec:bp-coefficient,b:sec:dyn-TV} were written for a general exponent $\lambda$. They use the deterministic assumptions stated there and the scale conditions \eqref{eq:ind-scales}, with $\rho_n=Rn^{-\lambda}$. \Cref{sec:references} is written for exponents satisfying \eqref{eq:ref-exponents} and an upper bound \eqref{eq:ref-mean-input} on the mean bias, and the static estimate \cref{b:lem:static-event} is stated for all scales satisfying the conditions \eqref{eq:st-S1}--\eqref{eq:st-S5}. In the proof below, the mean-bias input is the bound $0\le P_\beta-n^{-1}\E F_n\le C_bn^{-7/10}$ of \cref{inp:hb-rate}; the variance bound of \cref{inp:qp-rate} enters through \cref{sec:static}, and the trial bound of \cref{inp:qp-floor} through \cref{sec:fractional}.

\begin{proof}[Proof of \cref{thm:barrier-730}]
Fix $\beta>1$. We take the exponents $\lambda=7/30$ and $\xi=5/21$ in \eqref{eq:ref-exponents}, and the mean-bias input \eqref{eq:ref-mean-input} is the bound of \cref{inp:hb-rate}, with constant $C_b$. We fix the parameters of Steps 1--4 of \cref{tab:bp-order} as in \cref{sec:ref-parameters,sec:ref-families}, with $R^3>3C_b/(\varkappa g)$ as in \eqref{b:dyn:R-condition} and with an integer $B\ge H+2$ (for instance $B=H+2$). Then we fix $A_\beta$ as in Step 5, which is \eqref{b:dyn:coefficient} with $A=A_\beta$ and $E'=0$, and we set $L_n=A_\beta n^{7/30}$. The scales are $\rho_n=Rn^{-7/30}$, $\varepsilon_n=E\rho_n$, and $\eta_n=n^{-5/21}$. They satisfy \eqref{eq:ind-scales} with $E'=0$, since $\rho_n$, $\eta_n/\rho_n=n^{-1/210}/R$, and $1/(n\rho_n)=n^{-23/30}/R$ tend to zero.

We have $N_{\rm lev}\le t_0/(3\rho_n)=t_0n^{7/30}/(3R)\le e^{L_n}$ for all sufficiently large $n$. By \cref{rem:st-scales}(a), the scales satisfy \eqref{eq:st-S1}--\eqref{eq:st-S5}, so that \cref{b:lem:static-event} and therefore \cref{lem:ref-single-event} apply at these scales. Since $L_n=A_\beta n^{\lambda}$, the last assertion of \cref{lem:ref-families} shows that $L_n\to\infty$, $L_n/n\to0$, and \eqref{eq:ref-absorb} holds. By \cref{lem:ref-families}, applied with the time scale $L_n$, there exist measurable disorder events $\mathcal G_n$ with $\PP_J(\mathcal G_n)\to1$ such that, for all sufficiently large $n$ and every $J\in\mathcal G_n$, there exists a family of reference configurations with $N_{\rm lev}$ levels, with properties (i)--(iv) of that lemma, and with $\pi^J(V)\le3e^{-4L_n}$; as explained in \cref{sec:ind-setting}, it has properties (P1)--(P4). The family may depend on $J$ in an arbitrary way; it is not required to be a measurable function of $J$. Fix $\tau\in(0,1/2)$. By \cref{prop:bp-quenched}, for all sufficiently large $n$ and every $J\in\mathcal G_n$,
\[
 f_n(J)=\pi^J\bigl\{x:d^J(x,e^{L_n})\le\tau\bigr\}\le\frac{3(T_n+1)e^{-4L_n}+1/n}{1/2-\tau}.
\]
We write $e_n$ for the right side. The left side does not involve the family. By \cref{sec:model-joint}, $f_n$ is a Borel function of $J$, and $0\le f_n\le1$. By \eqref{eq:common-joint},
\[
 \PP_{J,X_0\sim\pi^J}\bigl\{d^J(X_0,e^{A_\beta n^{7/30}})\le\tau\bigr\}=\E_Jf_n(J)\le e_n+\PP_J(\mathcal G_n^c).
\]
Here $e_n\to0$ by \eqref{eq:bp-mixed-limit}, and $\PP_J(\mathcal G_n^c)\to0$. This proves \eqref{eq:barrier-730-continuous}. The same argument with \eqref{b:dyn:discrete-mixed} in place of \eqref{b:dyn:continuous-mixed} proves \eqref{eq:barrier-730-discrete}.

The coefficient $A_\beta$ was chosen in Step 5, before $\tau$, and the same coefficient serves both clocks. The threshold $\tau$ enters only through the factor $1/(1/2-\tau)$, which affects the rate of convergence and not the coefficient.
\end{proof}

The proof uses the two kinds of exceptional events of \cref{sec:static,sec:references} differently, as in \cref{rem:ref-two-errors}. The disorder probability $\PP_J(\mathcal G_n^c)$ is used once, additively, in the last display. It includes terms that decay only polynomially or as $e^{-cn\eta_n^4}$, and it is never multiplied by a number of references, levels, or time indices. At a fixed disorder in $\mathcal G_n$, the Gibbs mass $\pi^J(V)$ is multiplied by the $T_n+1$ time indices in \eqref{b:dyn:stationary-sign}, and the threshold probabilities are multiplied by the $T_n+1$ indices in the proof of \cref{prop:bp-obstruction}; both products are small.

\begin{proof}[Proof of \cref{b:cor:subcritical}]
Fix $\beta>1$, $\kappa\in(0,7/30)$, and $\tau\in(0,1/2)$, and let $A_\beta$ be the coefficient of \cref{thm:barrier-730}. Since $\kappa<7/30$, we have $n^\kappa\le A_\beta n^{7/30}$ for all sufficiently large $n$. By \cref{lem:model-dynamics}(d), the functions $t\mapsto d^J(x,t)$ and $k\mapsto d^J_{\rm disc}(x,k)$ are nonincreasing, and therefore, for all such $n$,
\[
 \bigl\{x:d^J(x,e^{A_\beta n^{7/30}})>\tau\bigr\}\subseteq\bigl\{x:d^J(x,e^{n^\kappa})>\tau\bigr\},
\]
and likewise in discrete time, with $\lfloor e^{n^\kappa}\rfloor\le\lfloor e^{A_\beta n^{7/30}}\rfloor$. By \eqref{eq:common-joint}, the realized-start probabilities are monotone in the same way, and \cref{thm:barrier-730} proves the corollary.
\end{proof}

\section{Quantitative and almost-sure forms}\label{sec:bp-strong}

The proof of \cref{thm:barrier-730} gives more than the convergence of the realized-start probabilities. In this section, we record four strengthenings. First, the disorder event can be chosen with a complement of stretched-exponentially small probability. Second, on this event there exists a set of initial states of Gibbs mass exponentially close to one from which the distance from equilibrium has a lower bound tending to $1/2$ up to the time horizon $e^{L_n}$. Third, the same holds for the Metropolis dynamics and for other single-site dynamics with bounded rates. Fourth, the relaxation time is at least of the order of the time horizon, up to a factor $n^{7/30}$. A consequence is the almost-sure form of the barrier along the full sequence of system sizes.

Three changes in the argument are needed. First, the lower bound on $F_n$ in the event of \cref{b:lem:static-event} is obtained from the small-deviation inequality of Paouris and Valettas for convex functions of Gaussian vectors \cite[Theorem~2.1 and Remark~2.4(3)]{paouris-valettas-2018}, in place of Chebyshev's inequality. Second, the tolerance $\eta_n$ is a fixed small multiple of $\rho_n$, so that the probability bounds for the fixed-scale constraints decay on the same scale $n\rho_n^4$ as the others. Third, the sign test of \cref{lem:bp-sign-test} is applied to each initial state separately, and the clock tail is bounded by a Chernoff bound instead of Markov's inequality. The relaxation-time bound uses in addition that the first reference takes few values. \Cref{thm:barrier-730} remains the main result; the coefficient in \cref{thm:barrier-strong} below is subject to a slightly smaller bound than in \eqref{b:dyn:coefficient} with $E'=0$, because of the second change.

\paragraph{Single-site dynamics.} Fix the disorder. For $x\in\Sigma_n$ and $1\le i\le n$, we write $x^{(i)}=x^{i,-x_i}$ for the configuration obtained by reversing the $i$th spin. For $C_{\rm rate}\ge1$, we let $\mathfrak L^J(C_{\rm rate})$ be the set of generators of the form
\begin{equation}\label{eq:bp-single-site}
 \mathcal Lf(x)=\sum_{i=1}^nc_i(x)\bigl(f(x^{(i)})-f(x)\bigr),\qquad f\colon\Sigma_n\to\R,
\end{equation}
with rates $c_i(x)\in[0,C_{\rm rate}]$, for which $\pi^J$ is stationary, that is, $\sum_x\pi^J(x)\mathcal Lf(x)=0$ for every $f$. We call $\mathcal L$ reversible if $\pi^J(x)c_i(x)=\pi^J(x^{(i)})c_i(x^{(i)})$ for all $x$ and $i$. By \eqref{eq:common-kernel}, the heat-bath generator $n(K^J-I)$ has the form \eqref{eq:bp-single-site} with $c_i(x)=p_i(-x_i\mid x)\le1$, and it is reversible by \cref{lem:model-dynamics}(a). The Metropolis kernel is
\[
 K_{\rm M}^J(x,x^{(i)})=\frac1n\min\bigl\{1,e^{\beta(H_n(x^{(i)})-H_n(x))}\bigr\},\qquad
 K_{\rm M}^J(x,x)=1-\sum_{i=1}^nK_{\rm M}^J(x,x^{(i)}),
\]
and $K_{\rm M}^J(x,y)=0$ otherwise. In one step, a uniformly chosen spin is reversed with the displayed acceptance probability. The identity $\pi^J(x)K^J_{\rm M}(x,x^{(i)})=(nZ_n)^{-1}\min\{e^{\beta H_n(x)},e^{\beta H_n(x^{(i)})}\}$, whose right side is symmetric in $x$ and $x^{(i)}$, shows that $K^J_{\rm M}$ is reversible with respect to $\pi^J$. Its continuous-time generator $n(K^J_{\rm M}-I)$ has the form \eqref{eq:bp-single-site} with $c_i(x)\le1$. Heat-bath and Metropolis dynamics belong to $\mathfrak L^J(1)$. For a reversible $\mathcal L$, the spectral gap $\operatorname{gap}(\mathcal L)$ is the infimum of $-\langle f,\mathcal Lf\rangle_{\pi^J}$ over $f$ with $\pi^J(f)=0$ and $\pi^J(f^2)=1$, where $\langle f,g\rangle_{\pi^J}=\sum_x\pi^J(x)f(x)g(x)$, and the relaxation time is $1/\operatorname{gap}(\mathcal L)$, with the convention $1/0=+\infty$.

\begin{theorem}[Quantitative and almost-sure forms]\label{thm:barrier-strong}
Fix $\beta>1$. There exist $A>0$, $c>0$, $C<\infty$, and Borel sets $\mathcal G_n^*$ of disorders, such that the following statements hold with $L_n=An^{7/30}$.
\begin{enumerate}[label=\textup{(\alph*)}]
\item For every $n\ge2$, we have $\PP_J\bigl((\mathcal G_n^*)^c\bigr)\le Ce^{-cn^{1/15}}$.
\item There exists $n_0$ such that, for every $n\ge n_0$ and every $J\in\mathcal G_n^*$, there exists a set $S\subseteq\Sigma_n$ with $\pi^J(S)\ge1-Ce^{-cn^{7/30}}$ such that every $x\in S$ satisfies
\[
 \inf_{0\le t\le e^{L_n}}d^J(x,t)\ge\frac12-Ce^{-cn^{7/30}},\qquad
 \min_{0\le k\le e^{L_n}}d^J_{\rm disc}(x,k)\ge\frac12-Ce^{-cn^{7/30}},
\]
and the same two inequalities hold for the Metropolis dynamics, that is, with $(K^J_{\rm M})^k$ in place of $(K^J)^k$ and $e^{nt(K^J_{\rm M}-I)}$ in place of $P^J_t$.
\item For every $C_{\rm rate}\ge1$ there exists $n_1$ such that, for every $n\ge n_1$, every $J\in\mathcal G_n^*$, and every $\mathcal L\in\mathfrak L^J(C_{\rm rate})$, there exists a set $S_{\mathcal L}\subseteq\Sigma_n$ with $\pi^J(S_{\mathcal L})\ge1-Ce^{-cn^{7/30}}$ such that
\[
 \norm{e^{t\mathcal L}(x,\cdot)-\pi^J}_\TV\ge\frac12-Ce^{-cn^{7/30}}\qquad\text{for every }x\in S_{\mathcal L}\text{ and }0\le t\le e^{L_n}.
\]
\item For every $C_{\rm rate}\ge1$, the threshold $n_1=n_1(\beta,C_{\rm rate})$ in part~(c) can be chosen so that, for every $n\ge n_1$, every $J\in\mathcal G_n^*$, and every reversible $\mathcal L\in\mathfrak L^J(C_{\rm rate})$, the relaxation time satisfies
\[
 \frac1{\operatorname{gap}(\mathcal L)}\ge\frac{e^{An^{7/30}}}{Cn^{7/30}}.
\]
In particular, this holds for the heat-bath and Metropolis generators $n(K^J-I)$ and $n(K^J_{\rm M}-I)$.
\item Let $(J^{(n)})_{n\ge2}$ be random disorders defined on one probability space, where $J^{(n)}$ has the law of the disorder of size $n$ and the dependence between different sizes is arbitrary. Then, almost surely, $J^{(n)}\in\mathcal G_n^*$ for all sufficiently large $n$. In particular, for every fixed $\tau\in(0,1/2)$, almost surely,
\[
 \pi^{J^{(n)}}\bigl\{x:d^{J^{(n)}}(x,e^{An^{7/30}})>\tau\bigr\}\longrightarrow1,\qquad
 \pi^{J^{(n)}}\bigl\{x:d^{J^{(n)}}_{\rm disc}(x,\lfloor e^{An^{7/30}}\rfloor)>\tau\bigr\}\longrightarrow1,
\]
and the same holds for the Metropolis dynamics.
\end{enumerate}
The constants $A$, $c$, $C$, and $n_0$ depend only on $\beta$; the threshold $n_1$ depends in addition on $C_{\rm rate}$.
\end{theorem}

Part~(b) contains the conclusions of \cref{thm:barrier-730}, with the coefficient $A$ of \cref{thm:barrier-strong} in place of $A_\beta$. It gives in addition the size of the exceptional sets, the uniformity over all earlier times, and a lower bound tending to $1/2$ on the distance. The threshold $1/2$ itself is not reached: the argument uses only that the half-space $\{y:R(v,y)\le0\}$ has Gibbs mass at least $1/2$. In part~(e), no independence between different system sizes is needed, because the first Borel--Cantelli lemma uses only the marginal laws.

The following table lists how the earlier statements enter the proof.

\begin{table}[ht]
\centering
\small
\begin{tabular}{@{}>{\raggedright\arraybackslash}p{0.36\textwidth}>{\raggedright\arraybackslash}p{0.56\textwidth}@{}}
\toprule
Statement & Use in the proof of \cref{thm:barrier-strong} \\
\midrule
\cref{sec:induction}; \cref{lem:bp-product}; \cref{prop:bp-obstruction} & used as stated; they assume \eqref{eq:ind-scales}, which holds with $E'=1/1000$ \\
\cref{lem:bp-coefficient} & used as stated, with $E'=1/1000$, so that \eqref{b:dyn:coefficient} reads $A<t_0\log(250/11)/(3BR)$; it covers every $m\le n^3e^{L_n}$ \\
\cref{b:lem:static-event} & statement used as stated, at $\eta_n=\rho_n/1000$; its hypotheses are verified in \cref{lem:bp-strong-event} \\
Step~1 of the proof of \cref{b:lem:static-event} (Chebyshev) & replaced by the Paouris--Valettas inequality in \cref{lem:bp-strong-event} \\
\cref{b:lem:coverage-extension} & used as stated, including its bound on $\PP(\mathcal C_n(D_1)^c)$ \\
\cref{lem:ref-single-event} & used as stated, at $\eta_n=\rho_n/1000$, in Step~4 of the proof of \cref{lem:bp-strong-event} \\
\cref{lem:ref-families} & used as stated, with $q_n$ supplied by that Step~4; \eqref{eq:ref-absorb} is verified there \\
\cref{lem:bp-sign-test}, \cref{prop:bp-quenched}, \eqref{b:dyn:clock-tail} & replaced by \cref{lem:bp-per-start} and the Chernoff bound \eqref{eq:bp-chernoff} \\
new & \cref{lem:bp-gap} (relaxation time) \\
\bottomrule
\end{tabular}
\caption{The earlier statements in the proof of \cref{thm:barrier-strong}.}
\label{tab:bp-strong-uses}
\end{table}

\begin{lemma}[A disorder event with stretched-exponential failure]\label{lem:bp-strong-event}
Fix $\beta>1$ and let $\lambda=7/30$. Fix the parameters of Steps 1--4 of \cref{tab:bp-order}, with $C_b$ the constant of \cref{inp:hb-rate}, and set
\[
 \rho_n=Rn^{-\lambda},\qquad\varepsilon_n=E\rho_n,\qquad\eta_n=\frac{\rho_n}{1000},\qquad L_n=An^\lambda
\]
for a fixed $A>0$. Let $\mathcal A_n$ be the event of \cref{b:lem:static-event} at these scales, with the data of \cref{sec:references}, and let $\mathcal G_n^*=\mathcal A_n\cap\bigcap_{s<B}\mathcal C_n(C_s+5)$ be the event \eqref{eq:ref-good-event} at these scales. Then the following hold.
\begin{enumerate}[label=\textup{(\roman*)}]
\item There exist $C,c>0$ such that $\PP_J((\mathcal G_n^*)^c)\le Ce^{-cn^{1-4\lambda}}$ for all sufficiently large $n$.
\item For all sufficiently large $n$ and every $J\in\mathcal G_n^*$, there exists a family of reference configurations with $N_{\rm lev}$ levels, with properties (P1)--(P4), and with $\pi^J(V)\le3e^{-4L_n}$.
\end{enumerate}
\end{lemma}

\begin{proof}
Since $\lambda>0$, we have $\rho_n=Rn^{-\lambda}\to0$. We also verify the conditions \eqref{eq:st-S1}--\eqref{eq:st-S5} of \cref{b:lem:static-event}. For \eqref{eq:st-S1}, we have $\Delta_{n,\beta}/\rho_n^3\le C_b/R^3<\varkappa g/3$, by \cref{inp:hb-rate} and \eqref{b:dyn:R-condition}. Next, $n^{7/15}/(n^2\rho_n^6)=R^{-6}n^{7/15-2+6\lambda}=R^{-6}n^{-2/15}$, which is \eqref{eq:st-S2}. We have $n\rho_n^4=R^4n^{1-4\lambda}=R^4n^{1/15}$ and $n\eta_n^4=10^{-12}R^4n^{1/15}$, which gives \eqref{eq:st-S3} and \eqref{eq:st-S4}. Finally, $\rho_n^3/\eta_n^2=10^6\rho_n\to0$, which is \eqref{eq:st-S5}. Therefore \cref{b:lem:static-event} applies at these scales.

\paragraph{Step 1: the numerators.} Steps 2--4 of the proof of \cref{b:lem:static-event} bound the probability that one of the events (A2)--(A4) fails by
\[
 C(n+1)^3\bigl(e^{-cn\rho_n^4}+e^{-cn\eta_n^4}\bigr)+C(n+1)e^{-cn}\le C'e^{-c'n^{1-4\lambda}},
\]
for all large $n$, where we used $n\eta_n^4=10^{-12}n\rho_n^4$ and $n\rho_n^4=R^4n^{1-4\lambda}$, and absorbed the polynomial factors by decreasing the constant in the exponent.

\paragraph{Step 2: the denominator.} The function $F_n=\log Z_n$ of the standard Gaussian vector $J=(J_{ij})_{i<j}$ is convex. Indeed, its Hessian is $\beta^2n^{-1}$ times the covariance matrix of the vector $(x_ix_j)_{i<j}$ under $\pi^J$, which is positive semidefinite. It is Lipschitz by \cref{rem:inp-lipschitz}, so it is square integrable. By \cref{inp:qp-rate}, $\sigma_n^2=\Var F_n\le C_vn^{7/15}$ for large $n$, with $C_v$ the variance constant there. We write $a_0<c_0$ for the static denominator margins in \eqref{b:eq:denominator-margin}. For all sufficiently large $n$, $\Delta_{n,\beta}\le a_0\rho_n^3$, so failure of (A1) implies $F_n-\E F_n<-u_n$, where $u_n=(c_0-a_0)n\rho_n^3$. If $\sigma_n=0$, this event is empty. Otherwise, let $t_n=u_n/\sigma_n$. Then
\[
 t_n\ge\frac{(c_0-a_0)R^3}{C_v^{1/2}}\,n^{23/30-3\lambda},
\]
and the exponent $23/30-3\lambda=1/15$ is positive, so that $t_n\ge2$ for large $n$. By \cite[Theorem~2.1 and Remark~2.4(3), (2.4)]{paouris-valettas-2018}, a convex square-integrable function $f$ of a standard Gaussian vector $Z$ satisfies
\begin{equation}\label{eq:bp-pv}
 \PP\Bigl(f(Z)-\E f(Z)<-t\sqrt{\Var f(Z)}\Bigr)\le\frac12\exp\Bigl(-\frac{\pi}{1024}(t-1)^2\Bigr)\qquad(t>1).
\end{equation}
% CHECK: arXiv:1611.01723v2, Theorem 2.1 (median form) and Remark 2.4(3), eq. (2.4) (mean form), p. 6.
Applying \eqref{eq:bp-pv} with $t=t_n$ and using $(t_n-1)^2\ge t_n^2/4$, we obtain
\[
 \PP\bigl(F_n<n(P_\beta-c_0\rho_n^3)\bigr)\le\frac12\exp\Bigl(-\frac{\pi t_n^2}{4096}\Bigr)
 \le\frac12\exp\Bigl(-\frac{\pi(c_0-a_0)^2R^6}{4096\,C_v}\,n^{23/15-6\lambda}\Bigr).
\]
Since $23/15-6\lambda=2/15$ is larger than $1-4\lambda=1/15$, this is at most $e^{-cn^{1-4\lambda}}$ for large $n$. Only the full partition function enters here. For the constrained numerators, the upper tail is needed, and the bounds of Step~1 remain those of \cref{b:lem:static-event}.

\paragraph{Step 3: the overlap-neighborhood events.} Fix $D_1>0$. By \cref{b:lem:coverage-extension}, there exists $D_4>0$, depending only on $\beta$ and $D_1$, such that, for all sufficiently large $n$,
\[
 \PP\bigl(\mathcal C_n(D_1)^c\bigr)\le3e^{-D_4L_n}=3e^{-D_4An^\lambda}.
\]
Since $\lambda=7/30$ is larger than $1-4\lambda=1/15$, this is at most $e^{-cn^{1-4\lambda}}$ for large $n$. There are $B$ such events in the definition of $\mathcal G_n^*$, with $D_1=C_s+5$ for $0\le s<B$. Combining Steps~1--3, we obtain part~(i).

\paragraph{Step 4: the families.} The present sequences satisfy $\rho_n\to0$ and \eqref{eq:st-S1}--\eqref{eq:st-S5}, as we checked at the start of the proof. By \cref{lem:ref-single-event}, for all large $n$ and every $J\in\mathcal A_n$, we have $(\pi^J)^{\otimes3}(\mathcal B_n)\le q_n$ with $q_n=C(n+1)^3(e^{-cn\rho_n^3}+e^{-cn\eta_n^2}+e^{-cn})$, where $\mathcal B_n$ is defined by \eqref{eq:ref-exceptional} with the present $\eta_n$. Since $n\rho_n^3=R^3n^{1-3\lambda}$ and $n\eta_n^2=10^{-6}R^2n^{1-2\lambda}$, and $1-3\lambda=3/10$ and $1-2\lambda=8/15$ are larger than $\lambda=7/30$, we have $e^{CL_n}q_n\to0$ for every fixed $C$, which is \eqref{eq:ref-absorb}. Moreover $L_n\to\infty$, $L_n/n\to0$, and $N_{\rm lev}\le t_0n^\lambda/(3R)\le e^{L_n}$ for large $n$. \Cref{lem:ref-families} therefore provides, for all large $n$ and every $J\in\mathcal G_n^*$, a family with properties (i)--(iv) of that lemma and $\pi^J(V)\le3e^{-4L_n}$. As in \cref{sec:ind-setting}, these imply (P1)--(P4), with $\zeta=2\eta_n$. This proves part~(ii).
\end{proof}

The next lemma refines \cref{lem:bp-sign-test}: it bounds the distance from equilibrium for each initial state separately, in terms of the probability that the chain started there visits $V$.

\begin{lemma}[The sign test for each initial state]\label{lem:bp-per-start}
Let $Q$, $\pi$, $v$, $V$, and $m$ satisfy the hypotheses of \cref{lem:bp-sign-test}. Let $(X_a)_{a\ge0}$ be the Markov chain with kernel $Q$, write $\PP_x$ for its law when $X_0=x$, and set
\[
 \phi(x)=\PP_x\{X_a\in V\text{ for some }0\le a\le m\},\qquad G_x=\{y:R(v(x),y)\le0\}.
\]
Then the following hold.
\begin{enumerate}[label=\textup{(\alph*)}]
\item $\sum_x\pi(x)\phi(x)\le(m+1)\pi(V)$, and hence $\pi\{x:\phi(x)>u\}\le(m+1)\pi(V)/u$ for every $u>0$.
\item For every $x$ and every $0\le k\le m$, we have $Q^k(x,G_x)\le \phi(x)$ and $\norm{Q^k(x,\cdot)-\pi}_\TV\ge1/2-\phi(x)$.
\item For every $x$, every $q_{\rm u}>0$, and every $t\ge0$, let $\mathsf N$ have the Poisson law with mean $q_{\rm u}t$. Then
\[
 \begin{gathered}
 \norm{e^{q_{\rm u}t(Q-I)}(x,\cdot)-\pi}_\TV\ge\frac12-\phi(x)-\PP\{\mathsf N>m\},\\
 \sum_x\pi(x)e^{q_{\rm u}t(Q-I)}(x,G_x)\le(m+1)\pi(V)+\PP\{\mathsf N>m\}.
 \end{gathered}
\]
\end{enumerate}
\end{lemma}

\begin{proof}
For part~(a), we have $\sum_x\pi(x)\phi(x)=\PP_\pi\{X_a\in V\text{ for some }a\le m\}$, where $\PP_\pi$ is the law of the chain with $X_0\sim\pi$. Since $\pi Q=\pi$, the law of $X_a$ under $\PP_\pi$ is $\pi$ for every deterministic $a$, and a union bound over $0\le a\le m$ gives the first assertion, as in \eqref{b:dyn:stationary-sign}. The second assertion is Markov's inequality. For part~(b), almost surely under $\PP_x$ we have $Q(X_{a-1},X_a)>0$ for every $a\ge1$. By the hypothesis of \cref{lem:bp-sign-test}, on the event that $X_0,\dots,X_k$ avoid $V$ we have $R(v(x),X_k)>0$, that is, $X_k\notin G_x$. It follows that $Q^k(x,G_x)\le\PP_x\{X_a\in V\text{ for some }a\le k\}\le \phi(x)$. As in the proof of \cref{lem:bp-sign-test}, spin-flip symmetry gives $\pi(G_x)\ge1/2$, and therefore $\norm{Q^k(x,\cdot)-\pi}_\TV\ge\pi(G_x)-Q^k(x,G_x)\ge1/2-\phi(x)$. For part~(c), the matrices $q_{\rm u}tQ$ and $q_{\rm u}tI$ commute, so $e^{q_{\rm u}t(Q-I)}=\sum_{k\ge0}\PP\{\mathsf N=k\}Q^k$. By part~(b),
\[
 e^{q_{\rm u}t(Q-I)}(x,G_x)\le\sum_{k\le m}\PP\{\mathsf N=k\}\phi(x)+\PP\{\mathsf N>m\}\le \phi(x)+\PP\{\mathsf N>m\},
\]
which gives the first inequality as in part~(b). Multiplying by $\pi(x)$, summing over $x$, and using part~(a), we obtain the second.
\end{proof}

The next lemma combines the standard continuous-time $L^2$ spectral-gap
contraction \cite[Lemma~20.5]{levin-peres-2017} with a partition of the
initial states according to their chosen reference. This grouping gives a
bound in terms of the number of references.

\begin{lemma}[Relaxation time and a sign test with few references]\label{lem:bp-gap}
Let $\mathcal L$ be a generator on $\Sigma_n$ that is reversible with respect to a probability measure $\pi$ with $\pi(-x)=\pi(x)$, and let $\operatorname{gap}(\mathcal L)$ be its spectral gap. Let $v\colon\Sigma_n\to\Sigma_n$ take at most $N_v$ values, and let $G_x=\{y:R(v(x),y)\le0\}$. If $t>0$ and
\[
 \omega=\sum_x\pi(x)e^{t\mathcal L}(x,G_x)<\frac12,
\]
then $\operatorname{gap}(\mathcal L)\,t\le(\log N_v)/2-\log(1-2\omega)$.
\end{lemma}

\begin{proof}
For $f$ with $\pi(f)=0$, the function $f_t=e^{t\mathcal L}f$ also has mean zero, since $\pi$ is stationary, and the derivative of $\langle f_t,f_t\rangle_\pi$ in $t$ is $2\langle f_t,\mathcal Lf_t\rangle_\pi\le-2\operatorname{gap}(\mathcal L)\langle f_t,f_t\rangle_\pi$, by the definition of the spectral gap. Gronwall's inequality gives $\norm{e^{t\mathcal L}f}_{L^2(\pi)}\le e^{-\operatorname{gap}(\mathcal L)t}\norm f_{L^2(\pi)}$.

For each value ${\widetilde v}$ of $v$, let $U_{\widetilde v}=\{x:v(x)={\widetilde v}\}$ and $H_{\widetilde v}=\{y:R({\widetilde v},y)\le0\}$. The sets $U_{\widetilde v}$ partition $\Sigma_n$, and $G_x=H_{\widetilde v}$ for $x\in U_{\widetilde v}$. As in the proof of \cref{lem:bp-sign-test}, spin-flip symmetry gives $\pi(H_{\widetilde v})\ge1/2$, and therefore
\[
 \sum_{\widetilde v}\pi(U_{\widetilde v})\pi(H_{\widetilde v})\ge\frac12,\qquad
 \sum_{\widetilde v}\langle\mathbf 1_{U_{\widetilde v}},e^{t\mathcal L}\mathbf 1_{H_{\widetilde v}}\rangle_\pi=\sum_x\pi(x)e^{t\mathcal L}(x,G_x)=\omega.
\]
Since $e^{t\mathcal L}\mathbf 1=\mathbf 1$ and $\pi$ is stationary, $\langle\mathbf 1_{U},e^{t\mathcal L}\mathbf 1_{U'}\rangle_\pi-\pi(U)\pi(U')=\langle\mathbf 1_U-\pi(U),e^{t\mathcal L}(\mathbf 1_{U'}-\pi(U'))\rangle_\pi$ for all sets $U,U'\subseteq\Sigma_n$. Using the Cauchy--Schwarz inequality, the contraction just proved, $\norm{\mathbf 1_U-\pi(U)}_{L^2(\pi)}^2=\pi(U)(1-\pi(U))\le\pi(U)$, and $\norm{\mathbf 1_{U'}-\pi(U')}_{L^2(\pi)}\le1/2$ for $U=U_{\widetilde v}$ and $U'=H_{\widetilde v}$, we find
\[
 \frac12-\omega\le\sum_{\widetilde v}\bigl(\pi(U_{\widetilde v})\pi(H_{\widetilde v})-\langle\mathbf 1_{U_{\widetilde v}},e^{t\mathcal L}\mathbf 1_{H_{\widetilde v}}\rangle_\pi\bigr)
 \le\frac{e^{-\operatorname{gap}(\mathcal L)t}}2\sum_{\widetilde v}\sqrt{\pi(U_{\widetilde v})}\le\frac{e^{-\operatorname{gap}(\mathcal L)t}}2\sqrt{N_v}.
\]
In the last step, we used the Cauchy--Schwarz inequality and $\sum_{\widetilde v}\pi(U_{\widetilde v})=1$. Rearranging gives the claim.
\end{proof}

The bound depends on the number $N_v$ of values of $v$, not on the size of $\Sigma_n$. If $v$ took all $2^n$ values, the lemma would only give $\operatorname{gap}(\mathcal L)\,t\le Cn$. The first reference takes at most $e^{CL_n}$ values, and this gives the factor $n^{7/30}$, rather than $n$, in \cref{thm:barrier-strong}(d).

\begin{proof}[Proof of \cref{thm:barrier-strong}]
Fix $\beta>1$ and let $\lambda=7/30$. We fix the parameters of Steps 1--4 of \cref{tab:bp-order} as in \cref{sec:ref-parameters,sec:ref-families}, with $C_b$ the constant of \cref{inp:hb-rate}, with $R^3>3C_b/(\varkappa g)$, and with an integer $B\ge H+2$. We set $E'=1/1000=E/10$, $\eta_n=\rho_n/1000$, and fix $A$ with
\[
 0<A<\frac{t_0\log(250/11)}{3BR},
\]
which is \eqref{b:dyn:coefficient} with $E'=1/1000$, since $4(E+E')=11/250$. Let $L_n=An^\lambda$, let $\mathcal G_n^*$ be the event of \cref{lem:bp-strong-event}, and let $m_*=\lfloor n^3e^{L_n}\rfloor$, a budget of update attempts. Throughout, $n$ is large enough that the conclusions of \cref{lem:bp-strong-event,lem:ind-margins,lem:bp-coefficient} hold; this fixes $n_0$.

\paragraph{Part (a).} By \cref{lem:bp-strong-event}(i) and $1-4\lambda=1/15$, the bound holds for $n\ge n_0$, and we enlarge $C$ so that $Ce^{-cn^{1/15}}\ge1$ for $2\le n<n_0$.

\paragraph{The obstruction.} Fix $n\ge n_0$ and $J\in\mathcal G_n^*$, and fix a family as in \cref{lem:bp-strong-event}(ii). The scales satisfy \eqref{eq:ind-scales} with $E'=1/1000$. By \cref{lem:ind-margins}, the margins \eqref{b:dyn:finite-margins} hold, and by \cref{lem:bp-coefficient}, $C_n(m)<1$ for every $m\le m_*$. By \cref{prop:bp-obstruction}, every path $x_0,\dots,x_m$ with $m\le m_*$ that avoids $V$ satisfies $R(v(x_0),x_m)>0$.

\paragraph{Parts (b) and (c).} We use the uniformization of a bounded-rate generator into a discrete kernel and a Poisson clock \cite[Section~20.1]{levin-peres-2017}. Let $\mathcal L\in\mathfrak L^J(C_{\rm rate})$ and set $Q_{\mathcal L}=I+\mathcal L/(C_{\rm rate}n)$. Its off-diagonal entries are $Q_{\mathcal L}(x,x^{(i)})=c_i(x)/(C_{\rm rate}n)\ge0$, its diagonal entries are $1-\sum_ic_i(x)/(C_{\rm rate}n)\ge0$, and $Q_{\mathcal L}(x,y)=0$ unless $y=x$ or $y=x^{(i)}$ for some $i$. We see that $Q_{\mathcal L}$ is a Markov kernel, $\pi^J$ is stationary for it, every sequence with $Q_{\mathcal L}(x_{a-1},x_a)>0$ is a path in the sense of \eqref{eq:ind-jump}, and $e^{t\mathcal L}=e^{C_{\rm rate}nt(Q_{\mathcal L}-I)}$. For the heat-bath and Metropolis generators with $C_{\rm rate}=1$, we have $Q_{\mathcal L}=K^J$ and $Q_{\mathcal L}=K^J_{\rm M}$, respectively. Since $\pi^J(-x)=\pi^J(x)$, the obstruction shows that the hypotheses of \cref{lem:bp-sign-test} hold for $Q=Q_{\mathcal L}$, $\pi=\pi^J$, the map $v$, the set $V$ of the family, and $m=m_*$. Let $\phi_{\mathcal L}$ be the function $\phi$ of \cref{lem:bp-per-start} for this kernel, and let
\[
 S_{\mathcal L}=\{x:\phi_{\mathcal L}(x)\le e^{-L_n}\}.
\]
Since $m_*+1\le2n^3e^{L_n}$ and $\pi^J(V)\le3e^{-4L_n}$, by \cref{lem:bp-per-start}(a), we have
\begin{equation}\label{eq:bp-strong-starts}
 \pi^J(S_{\mathcal L}^c)\le e^{L_n}(m_*+1)\pi^J(V)\le6n^3e^{-2L_n}.
\end{equation}
For $x\in S_{\mathcal L}$ and $0\le k\le e^{L_n}\le m_*$, by \cref{lem:bp-per-start}(b), we have $\norm{Q_{\mathcal L}^k(x,\cdot)-\pi^J}_\TV\ge1/2-e^{-L_n}$. For continuous time, the bound at $t=0$ is immediate. Let $0<t\le e^{L_n}$ and let $\mathsf N$ have the Poisson law with mean $\mu_{\rm P}=C_{\rm rate}nt\le C_{\rm rate}ne^{L_n}$. For every integer $j\ge\mu_{\rm P}$, Markov's inequality applied to $z^{\mathsf N}$ with $z=j/\mu_{\rm P}\ge1$, together with $\E z^{\mathsf N}=e^{\mu_{\rm P}(z-1)}$, gives $\PP\{\mathsf N\ge j\}\le z^{-j}e^{\mu_{\rm P}(z-1)}=e^{-\mu_{\rm P}}(e\mu_{\rm P}/j)^j$. With $j=m_*+1\ge n^3e^{L_n}$, this implies
\begin{equation}\label{eq:bp-chernoff}
 \PP\{\mathsf N>m_*\}\le\Bigl(\frac{e\,C_{\rm rate}ne^{L_n}}{n^3e^{L_n}}\Bigr)^{m_*+1}=\Bigl(\frac{eC_{\rm rate}}{n^2}\Bigr)^{m_*+1}\le e^{-(m_*+1)}\le e^{-L_n}
\end{equation}
once $n^2\ge e^2C_{\rm rate}$. By \cref{lem:bp-per-start}(c) with $q_{\rm u}=C_{\rm rate}n$, every $x\in S_{\mathcal L}$ satisfies $\norm{e^{t\mathcal L}(x,\cdot)-\pi^J}_\TV\ge1/2-2e^{-L_n}$ for $0\le t\le e^{L_n}$. Taking $n_1\ge n_0$ so large that also $n_1^2\ge e^2C_{\rm rate}$ and $12n^3e^{-2L_n}\le e^{-L_n}$ for $n\ge n_1$, we obtain part~(c) with $c=A$ and $C=2$. For part~(b), we apply this with $C_{\rm rate}=1$ to the heat-bath and Metropolis generators, for which $Q_{\mathcal L}$ is $K^J$ and $K^J_{\rm M}$, and take $S=S_{n(K^J-I)}\cap S_{n(K^J_{\rm M}-I)}$. By \eqref{eq:bp-strong-starts}, $\pi^J(S^c)\le12n^3e^{-2L_n}\le e^{-L_n}$, and the four distance bounds hold on $S$, with $\lfloor e^{L_n}\rfloor$ discrete steps and $1/2-2e^{-L_n}$ as the lower bound. Enlarging $n_0$ to the value of $n_1$ for $C_{\rm rate}=1$ completes part~(b).

\paragraph{Part (d).} Let $\mathcal L\in\mathfrak L^J(C_{\rm rate})$ be reversible, and let $t=e^{L_n}$. By \cref{lem:bp-per-start}(c) with $q_{\rm u}=C_{\rm rate}n$, by \eqref{eq:bp-chernoff}, and by $\pi^J(V)\le3e^{-4L_n}$,
\[
 \omega=\sum_x\pi^J(x)e^{t\mathcal L}(x,G_x)\le6n^3e^{-3L_n}+e^{-L_n}\le\frac14
\]
for $n\ge n_1$, after enlarging the threshold $n_1$ of part~(c) by an amount depending only on $\beta$ and $C_{\rm rate}$. Write $\widehat c=c(5)>0$ for the first list-size constant in \eqref{eq:ref-schedule}. The map $v$ takes values among the signed entries of the list $\mathcal F_0$, which has $K_0=\lceil e^{(\widehat c+1)L_n}\rceil$ entries by \eqref{eq:ref-schedule}; hence it takes at most $N_v=2K_0\le4e^{(\widehat c+1)L_n}$ values, where $\widehat c$ depends only on $\beta$. By \cref{lem:bp-gap},
\[
 \operatorname{gap}(\mathcal L)e^{L_n}\le\frac12\log\bigl(4e^{(\widehat c+1)L_n}\bigr)+\log2\le(\widehat c+2)L_n
\]
for $n\ge n_1$, after enlarging $n_1$ once more by an amount depending only on $\beta$. If $\operatorname{gap}(\mathcal L)=0$, the conclusion of part~(d) is immediate by the convention $1/0=+\infty$. Otherwise we may divide by the positive gap, and we obtain $1/\operatorname{gap}(\mathcal L)\ge e^{L_n}/((\widehat c+2)L_n)=e^{An^\lambda}/((\widehat c+2)An^\lambda)$, which is part~(d) after enlarging $C$. The enlarged $n_1$ also serves in part~(c).

\paragraph{Part (e).} By part~(a), $\sum_n\PP(J^{(n)}\notin\mathcal G_n^*)\le\sum_nCe^{-cn^{1/15}}<\infty$. By the first Borel--Cantelli lemma, which uses only the marginal laws of the $J^{(n)}$, almost surely $J^{(n)}\in\mathcal G_n^*$ for all large $n$. On this event, part~(b) gives, for all large $n$, a set of Gibbs mass at least $1-Ce^{-cn^\lambda}$ on which $d^{J^{(n)}}(x,e^{An^\lambda})$ and $d^{J^{(n)}}_{\rm disc}(x,\lfloor e^{An^\lambda}\rfloor)$ are at least $1/2-Ce^{-cn^\lambda}$, which exceeds $\tau$ for large $n$, and likewise for the Metropolis dynamics. This completes the proof.
\end{proof}
